\section{Design}
\label{sec:method-design}

The design is the simplest one that answers the research questions: a single training
pipeline with exactly two pluggable variables, and one run per cell of the resulting
grid. Table~\ref{tab:factorial} gives the twelve runs.

\begin{table}[htbp]
\centering
\caption{The twelve experiments. E1 to E6 are the confirmatory factorial, specified
before any results were seen. E7 to E12 are the EfficientNet-B3 ablation, added
afterwards and analysed separately.}
\label{tab:factorial}
\small
\begin{tabular}{@{}llccc@{}}
\toprule
\textbf{Architecture} & \textbf{Input} & \textbf{Augmentation only} &
\textbf{Weighted CE} & \textbf{Oversampling} \\
\midrule
\multicolumn{5}{@{}l}{\emph{Confirmatory factorial}} \\
MobileNetV2     & 224 & E1 & E3 & E5 \\
ResNet-50       & 224 & E2 & E4 & E6 \\
\midrule
\multicolumn{5}{@{}l}{\emph{Ablation}} \\
EfficientNet-B3 & 224 & E7  & E8  & E9  \\
EfficientNet-B3 & 300 & E10 & E11 & E12 \\
\bottomrule
\end{tabular}
\end{table}

The experiment registry is the single place where these definitions live. Nothing
downstream hard-codes an architecture or a resolution; every script and notebook reads
from this dictionary, which is what keeps the twelve runs from drifting apart.

\begin{lstlisting}[caption={The experiment registry, \texttt{src/ham10000/constants.py}.},
                   label={lst:registry}]
EXPERIMENTS = {
    # --- confirmatory 2x3 factorial, 224px ---
    "E1":  {"architecture": "mobilenet_v2",    "strategy": "augmentation", "image_size": 224},
    "E2":  {"architecture": "resnet50",        "strategy": "augmentation", "image_size": 224},
    "E3":  {"architecture": "mobilenet_v2",    "strategy": "weighted_ce",  "image_size": 224},
    "E4":  {"architecture": "resnet50",        "strategy": "weighted_ce",  "image_size": 224},
    "E5":  {"architecture": "mobilenet_v2",    "strategy": "oversampling", "image_size": 224},
    "E6":  {"architecture": "resnet50",        "strategy": "oversampling", "image_size": 224},
    # --- EfficientNet-B3 ablation, matched 224px ---
    "E7":  {"architecture": "efficientnet_b3", "strategy": "augmentation", "image_size": 224},
    "E8":  {"architecture": "efficientnet_b3", "strategy": "weighted_ce",  "image_size": 224},
    "E9":  {"architecture": "efficientnet_b3", "strategy": "oversampling", "image_size": 224},
    # --- EfficientNet-B3 ablation, native 300px ---
    "E10": {"architecture": "efficientnet_b3", "strategy": "augmentation", "image_size": 300},
    "E11": {"architecture": "efficientnet_b3", "strategy": "weighted_ce",  "image_size": 300},
    "E12": {"architecture": "efficientnet_b3", "strategy": "oversampling", "image_size": 300},
}
\end{lstlisting}

\section{What Varies and What Does Not}
\label{sec:method-contract}

Table~\ref{tab:contract} is the experimental contract. Everything in the left column is
identical across all twelve runs; only the right column changes.

\begin{table}[htbp]
\centering
\caption{Experimental contract.}
\label{tab:contract}
\small
\renewcommand{\arraystretch}{1.2}
\begin{tabular}{@{}p{7.4cm}p{5.6cm}@{}}
\toprule
\textbf{Held identical across all twelve} & \textbf{Varied} \\
\midrule
Split: lesion-level, stratified 80/10/10, seed 42, hash \texttt{f35ac1de18678182}
  & \textbf{Architecture}: MobileNetV2, ResNet-50, EfficientNet-B3 \\
Augmentation: HFlip, VFlip, rotation $\pm 30^{\circ}$, colour jitter 0.2
  & \textbf{Imbalance strategy}: augmentation-only, weighted CE, oversampling \\
Normalisation: ImageNet channel statistics
  & \textbf{Input resolution}: 224 or 300 (EfficientNet-B3 only) \\
Initialisation: ImageNet-pretrained, classifier head reset to 7 outputs & \\
Optimiser: Adam, lr $=10^{-3}$, $\beta = (0.9, 0.999)$ & \\
Batch size: 32 & \\
Schedule: max 50 epochs, early stopping on validation loss, patience 10 & \\
Dataloader workers: 16 & \\
Seed: 42 & \\
Hardware: one NVIDIA RTX A4500 for every run & \\
\bottomrule
\end{tabular}
\end{table}

Two entries in that table are less obvious than the rest and are worth justifying.

\paragraph{Dataloader workers belong in the contract.} PyTorch seeds each dataloader
worker from \texttt{base\_seed + worker\_id}, and the augmentation transforms draw from
that per-worker generator. Changing the worker count therefore changes which random flip
and rotation each image receives. It biases no condition, but it must be constant across
runs for the comparison to remain controlled, which is why it is listed here rather than
treated as an implementation detail.

\paragraph{Batch size is fixed even where memory would allow more.} The 300-pixel runs
have headroom on the GPU used. Increasing their batch size would speed them up and would
also make them incomparable with the other nine, so it was not done.

\section{Architectures}
\label{sec:method-arch}

All three backbones are loaded from \texttt{torchvision} with ImageNet-pretrained
weights. Only the final linear layer is replaced, mapping to seven outputs. No layers are
frozen; every run fine-tunes end to end.

\begin{lstlisting}[caption={Backbone construction, \texttt{src/ham10000/models.py}. For
MobileNetV2 and EfficientNet-B3 the classifier is a \texttt{Sequential(Dropout, Linear)},
so replacing index 1 leaves the dropout at index 0 in place.},
                   label={lst:models}]
def build_model(architecture: str, num_classes: int = 7) -> nn.Module:
    arch = architecture.lower()

    if arch == "mobilenet_v2":
        model = models.mobilenet_v2(weights=models.MobileNet_V2_Weights.IMAGENET1K_V1)
        model.classifier[1] = nn.Linear(model.classifier[1].in_features, num_classes)
        return model

    if arch in ("resnet50", "resnet_50"):
        model = models.resnet50(weights=models.ResNet50_Weights.IMAGENET1K_V1)
        model.fc = nn.Linear(model.fc.in_features, num_classes)
        return model

    if arch in ("efficientnet_b3", "efficientnet-b3"):
        model = models.efficientnet_b3(weights=models.EfficientNet_B3_Weights.IMAGENET1K_V1)
        model.classifier[1] = nn.Linear(model.classifier[1].in_features, num_classes)
        return model

    raise ValueError(f"Unknown architecture: {architecture}. Expected one of {SUPPORTED}.")
\end{lstlisting}

Table~\ref{tab:arch-structure} reports the structural properties of the three networks
as instantiated, counted from the constructed modules rather than quoted from the
original papers.

\begin{table}[htbp]
\centering
\caption{Structural properties, counted from the built models with a seven-class head.}
\label{tab:arch-structure}
\small
\begin{tabular}{@{}lrrrl@{}}
\toprule
\textbf{Architecture} & \textbf{Trainable params} & \textbf{Weight layers} &
\textbf{Blocks} & \textbf{Dropout} \\
\midrule
MobileNetV2     &  2{,}232{,}839 &  53 & 17 inverted residual & $p = 0.2$ \\
ResNet-50       & 23{,}522{,}375 &  54 & 16 bottleneck        & none \\
EfficientNet-B3 & 10{,}706{,}991 & 131 & 26 MBConv            & $p = 0.3$ \\
\bottomrule
\end{tabular}
\end{table}

Two observations from that table shape later interpretation.

First, ResNet-50's 54 weight layers exceed the 50 in its name. The
\texttt{torchvision} implementation contains 53 convolutional modules because four are
$1\times1$ projection convolutions on the shortcut path at stage transitions, which the
canonical count omits. The figure is reported as counted so that it can be reconciled
with the code.

Second, EfficientNet-B3 is by far the deepest at 131 weight layers while holding less
than half of ResNet-50's parameters. It is a deep, narrow network where ResNet-50 is
shallower and wider. Chapter~7 shows this has consequences for inference latency that
parameter count alone does not predict.

\subsection{An asymmetry in dropout}
\label{sec:method-dropout}

The three backbones do not carry the same regularisation. MobileNetV2 ships with a
dropout layer at $p = 0.2$ before its classifier, EfficientNet-B3 with $p = 0.3$, and
ResNet-50 with none at all. Because only the final \texttt{Linear} is replaced, each
network retains whatever its reference implementation provides.

This is left as-is deliberately, and it is a defensible choice rather than an oversight,
but it is also a genuine asymmetry between the arms and is declared as such. The
comparison is intended to be between these networks as they are published and as anyone
would actually deploy them. Matching dropout across all three would compare three
variants that exist nowhere else. The cost is that a small part of any advantage held by
the two networks with dropout could be attributable to the regularisation rather than to
the architecture. Chapter~9 records the matched-dropout ablation as future work; it is a
one-line change.

\section{Imbalance Strategies}
\label{sec:method-strategies}

\begin{lstlisting}[caption={The three strategies, \texttt{src/ham10000/imbalance.py}.
Each returns a dataloader and the loss function that goes with it.},
                   label={lst:imbalance}]
def make_train_loader(dataset, train_labels, strategy, batch_size=32,
                      num_workers=2, pin_memory=True):
    strategy = strategy.lower()
    if strategy == "augmentation":
        loader = DataLoader(dataset, batch_size=batch_size, shuffle=True,
                            num_workers=num_workers, pin_memory=pin_memory)
        return loader, torch.nn.CrossEntropyLoss()

    if strategy == "weighted_ce":
        class_weights = inverse_frequency_weights(train_labels)
        loader = DataLoader(dataset, batch_size=batch_size, shuffle=True,
                            num_workers=num_workers, pin_memory=pin_memory)
        return loader, torch.nn.CrossEntropyLoss(weight=class_weights)

    if strategy == "oversampling":
        weights = sample_weights_for_oversampling(train_labels)
        sampler = WeightedRandomSampler(weights, num_samples=len(weights), replacement=True)
        loader = DataLoader(dataset, batch_size=batch_size, sampler=sampler,
                            num_workers=num_workers, pin_memory=pin_memory)
        return loader, torch.nn.CrossEntropyLoss()

    raise ValueError(f"Unknown imbalance strategy: {strategy}")
\end{lstlisting}

\paragraph{Augmentation-only (E1, E2, E7, E10).} Standard cross-entropy with shuffled
sampling. The augmentation pipeline of Section~\ref{sec:data-preproc} still applies, as
it does everywhere, but nothing is done about imbalance. This is the no-remediation
control against which the other two are read.

\paragraph{Weighted cross-entropy (E3, E4, E8, E11).} Per-class weights from
Equation~\ref{eq:weights}, passed to the loss. Sampling is unchanged. The model sees the
same images in the same proportions; the gradient contribution of each is rescaled.

\paragraph{Random oversampling (E5, E6, E9, E12).} A
\texttt{WeightedRandomSampler} draws with replacement at probability proportional to
$1/n_c$, so an epoch contains roughly balanced class frequencies while the loss stays
unweighted. Rare-class images recur many times per epoch; common-class images are
frequently skipped. This is the condition that tests Buda et
al.'s~\cite{buda2018classimbalance} finding that oversampling does not induce
overfitting in CNNs.

The distinction between the last two is worth keeping clear, because they are often
treated as interchangeable. Weighted cross-entropy changes how much each example counts.
Oversampling changes how often each example is seen. They arrive at a similar objective
by different routes and, as Chapter~7 shows, do not produce the same result.

\section{Training}
\label{sec:method-training}

\begin{lstlisting}[caption={Optimiser construction, \texttt{src/ham10000/train.py}.},
                   label={lst:optim}]
# No weight decay, no LR schedule, no gradient clipping. Held identical
# across all twelve runs, and disclosed in the methodology chapter.
optimizer = torch.optim.Adam(model.parameters(), lr=config.lr, betas=(0.9, 0.999))
\end{lstlisting}

Adam at $10^{-3}$ with default $\beta$ values, and nothing else. That is a deliberate
omission and needs stating explicitly, because the absence of these is as much a design
decision as their presence would be.

\subsection{Regularisation applied, and not applied}
\label{sec:method-regularisation}

Four mechanisms restrain overfitting: the augmentation pipeline, early stopping on
validation loss with patience 10, ImageNet pre-training as a strong prior on the
weights, and the architecture-specific dropout of
Section~\ref{sec:method-dropout}.

What is not applied: no weight decay, since \texttt{weight\_decay} retains its default
of zero and there is no $L_2$ penalty; no learning-rate schedule; no gradient clipping;
no label smoothing. These omissions are uniform, so they confound no comparison between
conditions. They remain a real limitation, and Chapter~8 argues they are the most
plausible explanation for the optimisation behaviour observed in the ResNet-50 runs.

\subsection{Early stopping and the epoch budget}

Training runs to a maximum of 50 epochs with early stopping on validation loss at
patience 10. The best checkpoint by validation loss is the one evaluated on test.

The patience value was revised upward from an initial 5 during piloting, after inspecting
the validation curves rather than any test-set number. ResNet-50 was still finding new
best losses at the point where patience 5 terminated it, meaning the original budget was
cutting off a run that had not converged. The revision was applied uniformly to all
twelve conditions, so the comparison remains controlled. It is nonetheless a deviation
from the originally registered protocol and is disclosed here rather than absorbed
silently.

\subsection{Resumability and configuration drift}

Every epoch writes a checkpoint. A run interrupted by a dropped connection resumes from
the last completed epoch rather than restarting, which matters when twelve runs share
one rented GPU session.

Resumption introduces a hazard that is easy to miss: a checkpoint written under one
configuration can be resumed under another, splicing two protocols into a single history
file with nothing in the output to indicate it happened. The training loop therefore
compares the stored configuration against the current one and refuses to continue on any
mismatch.

\begin{lstlisting}[caption={Configuration drift guard, \texttt{src/ham10000/train.py}.},
                   label={lst:drift}]
_MUST_MATCH = ("architecture", "strategy", "image_size", "batch_size",
               "lr", "num_workers", "seed", "split_hash")
old = ckpt.get("config", {})
drift = {
    k: (old.get(k), run_config.get(k))
    for k in _MUST_MATCH
    if k in old and old.get(k) != run_config.get(k)
}
if drift:
    raise RuntimeError(
        f"{config.experiment_id} has a checkpoint from a different configuration:\n"
        f"Resuming would mix both protocols in one run."
    )
\end{lstlisting}

\section{Compute Environment}
\label{sec:method-compute}

All twelve runs execute on one rented NVIDIA RTX A4500 (20{,}470~MiB) under PyTorch
2.8.0 with CUDA 12.8 and torchvision 0.23.0~\cite{paszke2019pytorch}. Using a single GPU
for every run removes hardware as a source of variation: cuDNN kernel selection is not
identical across GPU architectures even with deterministic mode enabled, so runs split
across different cards would not be strictly comparable.

Peak GPU memory was measured for a real training step in each configuration before the
runs began, to confirm the card was adequate rather than assume it.

\begin{table}[htbp]
\centering
\caption{Peak GPU memory for a training step at batch size 32, measured on the RTX A4500.}
\label{tab:vram}
\small
\begin{tabular}{@{}lrl@{}}
\toprule
\textbf{Configuration} & \textbf{Peak} & \textbf{Experiments} \\
\midrule
MobileNetV2 @224     &  2.6 GB & E1, E3, E5 \\
ResNet-50 @224       &  3.1 GB & E2, E4, E6 \\
EfficientNet-B3 @224 &  5.8 GB & E7, E8, E9 \\
EfficientNet-B3 @300 & 10.2 GB & E10, E11, E12 \\
\bottomrule
\end{tabular}
\end{table}

The 300-pixel runs set the ceiling at 10.2~GB, comfortably inside the card. The dataset
was staged into a RAM-backed filesystem for training after profiling showed that reading
individual JPEGs from the network-attached volume, not GPU computation, was the limiting
factor: 45 images per second from the network mount against 241 from RAM. This affects
throughput only. The bytes, their order and the seeds are unchanged, so it has no bearing
on any reported result.

\section{Reproducibility}
\label{sec:method-repro}

Seeds are fixed at 42 for Python, NumPy and PyTorch, and cuDNN runs in deterministic
mode with benchmarking disabled.

\begin{lstlisting}[caption={Seeding, \texttt{src/ham10000/train.py}.}, label={lst:seed}]
def set_seed(seed: int) -> None:
    random.seed(seed)
    np.random.seed(seed)
    torch.manual_seed(seed)
    torch.cuda.manual_seed_all(seed)
    torch.backends.cudnn.deterministic = True
    torch.backends.cudnn.benchmark = False
\end{lstlisting}

Each run persists five artefacts: \texttt{config.json} with the fully resolved
configuration including the split hash, GPU name and library versions;
\texttt{history.csv} with per-epoch training and validation curves;
\texttt{test\_metrics.json} with the full metric set and confusion matrix;
\texttt{test\_predictions.csv} with a per-image prediction, required for the paired
statistical tests; and the best checkpoint by validation loss.

Every table and figure in Chapter~7 is generated from those artefacts by script rather
than transcribed. A number in this document that disagreed with the stored results would
require the generating code to be wrong, not merely a typo, which is a stronger guarantee
than proofreading provides.
